\documentclass{jreport}
\usepackage{amsmath}
\usepackage{amsfonts}
\usepackage{amsthm}
\usepackage{xcolor}
\usepackage{bm}
\usepackage[dvipdfmx,hidelinks]{hyperref}
\usepackage{pxjahyper}
\usepackage{amssymb}
\begin{document}
\title{第一回線形代数学レポート問題}
\author{鄒暁城 \\ C1SB1018}
\maketitle
\newpage
\setcounter{chapter}{9}
\tableofcontents
\addcontentsline{toc}{chapter}{目次}
\newpage
\section{次の行列の固有値と固有ベクトルを求めよ.}
$$
	A =
	\left(
	\begin{array}{cc}
		1&2\\
		5&4\\
	\end{array}
	\right)
,
	B=
	\left(
	\begin{array}{ccc}
		1&2&4\\
		2&-2&2\\
		4&2&1\\
	\end{array}
	\right)
$$
まず$A$に対し,固有多項式$\phi_A (t) =(1-t)(4-t)-10=(t+1)(t-6)$である.\\
よって,固有値は以下のように与えられる:
$$
	\begin{cases}
	\lambda_1 = -1\\
	\lambda_2 =  6
	\end{cases}
$$
簡単の計算より$\lambda_1$に属する固有ベクトル$\vec{v_1}$,$\lambda_2$に属する固有ベクトル$\vec{v_2}$は次のように定める:
$$
\vec{v_1}=
\left(
\begin{array}{c}
1\\
-1\\
\end{array}
\right)
,
\vec{v_2}=
\left(
\begin{array}{c}
2\\
5\\
\end{array}
\right)
$$
次に行列$B$に対し,固有多項式$\phi_B (t) =(t+3)^2(t-6)$である.\\
よって,固有値は次のように与えられる:
$$
        \begin{cases}
        \lambda_1 = -3\\
	\lambda_2 = -3\\
        \lambda_3 =  6
        \end{cases}
$$
簡単の計算より$\lambda_1$に属する固有ベクトル$\vec{v_1},\vec{v_2}$,$\lambda_3$に属する固有ベクトル$\vec{v_3}$は次のように定める:
$$
\vec{v_1}=
\left(
\begin{array}{c}
-1\\
0\\
1\\
\end{array}
\right)
,
\vec{v_2}=
\left(
\begin{array}{c}
-1\\
2\\
0\\
\end{array}
\right)
,
\vec{v_3}=
\left(
\begin{array}{c}
2\\
1\\
2\\
\end{array}
\right)
$$
\section{$n$次正方行列$A$について,$A$が正則行列である必要十分条件は,$A$が$0$を固有値として持たないことである}
$n$次正方行列$A$に対し,$A$が正則行列であることと$A\vec{x}=\vec{0}$が非自明な解をもたないことと同値である.\\
また,$n$次正方行列$A$が$0$を固有値に持たないことと,$A\vec{x}=\vec{0}$が非自明な解をもたないと同値であるため.問題の主張が自明である。
\section{次の行列の対角化可能性を判定し,可能の場合は対角化せよ}
$$
A=
\left(
\begin{array}{ccc}
	1 & 5 & 1\\
	4 & 0 & 1\\
	0 & 0 &-4\\
\end{array}
\right)
,
B=
\left(
\begin{array}{ccc}
        0 & 1 &-1\\
       -2 & 3 &-1\\
       -1 & 1 & 1\\
\end{array}
\right)
,
C=
\left(
\begin{array}{ccc}
        0 & 1 & 1\\
        1 & 0 & 1\\
        1 & 1 & 0\\
\end{array}
\right)
$$
行列$A$が対角化不可能である.固有多項式$\phi_A(t)=-(t+4)^2(t-5)$より,\\
固有値$\lambda_1 =-4$の重複度は$2$で,\\
しかし,
$$
V(\lambda_1) =Span[\left(
			\begin{array}{c}
				1\\
				-1\\
				0\\
			\end{array}
			\right)
			]
$$
上記の式より,$\dim(V(\lambda_1)) =1 <2$が成り立つため、対角化不可能である.\\
行列$B$も対角化不可能である.固有多項式$\phi_B(t)= (t-1)^2(t-2)$より,\\
固有値$\lambda_1 =1$の重複度が$2$で,\\
しかし,
$$
V(\lambda_1) =Span[\left(
                        \begin{array}{c}
                                1\\
                                 1\\
                                0\\
                        \end{array}
                        \right)
                        ]
$$
上記の式より,$\dim(V(\lambda_1)) =1 <2$が成り立つため、対角化不可能である.\\
行列$C$が対角か可能である.固有多項式$\phi_C(t) =(t+1)^2(2-t)$より,\\
固有値$\lambda_1 =-1$の重複度が$2$で,$\lambda_2 =2$の重複度が$1$であり、\\
また,
$$
V(\lambda_1) =Span[\left(
                        \begin{array}{c}
                                -1\\
                                1\\
                                0\\
                        \end{array}
                        \right)
			,
			\left(
			\begin{array}{c}
				-1\\
				0\\
				1\\
			\end{array}
			\right)
                        ]
$$
$$
V(\lambda_2) = Span[\left(
			\begin{array}{c}
				1\\
				1\\
				1\\
			\end{array}
			\right)
			]
$$
より,行列$P$を次のようにおくと:
$$
\left(
\begin{array}{ccc}
	-1 & -1 &  1\\
	 1 &  0 &  1\\
	 0 &  1 &  1\\
\end{array}
\right)
$$
次のような対角化ができる:\\
$$
P^{-1}CP=
\left(
\begin{array}{ccc}
	-1 & 0 & 0\\
	0  & -1 & 0\\
	0  &  0 & 2\\
\end{array}
\right)
$$
\section{対角化可能性を判定し,可能の場合対角化せよ}
$$
A=
\left(
\begin{array}{cccc}
	0 & 0 & 0 & 1\\
	0 & 0 & 1 & 0\\
	0 & 1 & 0 & 0\\
	1 & 0 & 0 & 0\\
\end{array}
\right)
,
B=
\left(
\begin{array}{ccc}
        2 & 2 & 2\\
        0 & 3 & 1\\
       -1 & 2 & 3\\
\end{array}
\right)
$$
行列$A$が対角化可能である.固有多項式$\phi_A(t) =(t-1)^2(t+1)^2$であるため。\\
固有値$\lambda_1 =1$の重複度が$2$で,固有値$\lambda_2 =-1$の重複度も$2$である.\\
この時行列$P$を次のように定めると,対角化ができる:
$$
P=
\left(
\begin{array}{cccc}
        0 & 1 & 0 &-1\\
        1 & 0 &-1 & 0\\
        1 & 0 & 1 & 0\\
        0 & 1 & 0 & 1\\
\end{array}
\right)
$$
以下のように対角化ができる:
$$
P^{-1}AP=
\left(
\begin{array}{cccc}
        1 & 0 & 0 & 0\\
        0 & 1 & 0 & 0\\
        0 & 0 &-1 & 0\\
        0 & 0 & 0 & -1\
\end{array}
\right)
$$
行列$B$は対角化不可能である,固有多項式$\phi_B(t) =(2-t)(2-t)^2$であるため.\\
固有値$\lambda_1 =2$の重複度が$1$で,固有値$\lambda_2 =3$の重複度が$1$である.\\
しかし
$$
V(\lambda_2) = Span[\left(
                        \begin{array}{c}
                                2\\
                                1\\
                                0\\
                        \end{array}
                        \right)
                        ]
$$
上記の式が成り立つことより,$\dim(V(\lambda_2)) =1 <2$より、対角化不可能である.
\section{対角化可能性を判定し,できるやつを対角化せよ}
$$
A=
\left(
\begin{array}{ccc}
        a & c & 0\\
        0 & a & 0\\
        0 & 0 & b\\
\end{array}
\right)
,
B=
\left(
\begin{array}{ccc}
        a & 0 & 0\\
        0 & a & c\\
        0 & 0 & b\\
\end{array}
\right)
$$
この時,2つの行列の固有多項式はいずれも$\phi (t) =(a-t)^2(b-t)$であり、\\
固有値$\lambda_1 = a$の重複度が$2$で,固有値$\lambda_2 =b$の重複度は$1$である.\\
このもとで,計算して議論すると、行列$A$が対角化不可能で、行列$B$が対角か可能である.\\
行列$A$に対し,$a \ne b,c\ne 0$を念頭において、\\
$$
V_A(\lambda_1)=Span[\left(
                        \begin{array}{c}
                                1\\
                                0\\
                                0\\
                        \end{array}
                        \right)
                        ]
$$
以上の式が成り立つことより,$\dim(V_A(\lambda_1)) =1 <2$となり、対角化が不可能である.\\
一方,行列$P$を以下のように定めると、行列$B$を対角化できる:
$$
P=
\left(
\begin{array}{ccc}
         1 &  0 &  0\\
         0 &  1 &  c\\
         0 &  0 &  b-a\\
\end{array}
\right)
$$
対角化を次のようにできる:
$$
P^{-1}BP=
\left(
\begin{array}{ccc}
         a &  0 &  0\\
         0 &  a &  0\\
         0 &  0 &  b\\
\end{array}
\right)
$$
\section{$n$次正方行列$A$の固有多項式は$n$重解を持つ時,単位行列の定数倍でないならば、対角化不可であることを示せ}
以下対偶を示す:\\
行列$A$が対角化可能であると仮定して,$A=\lambda E$であることを示す:\\
対角化可能と固有方程式の解の構造より:\\
ある正則行列$P$が存在し,$P^{-1}AP = \lambda E$が成り立つ.\\
これを変形して,$A=P\cdot \lambda E \cdot P^{-1} = \lambda PP^{-1} = \lambda E$となり、\\
対偶命題を示せた.\\
したがって,問題文の主張が正しいことを示せた.
\end{document}

